\documentclass{beamer}

\usepackage[utf8]{inputenc}
\input{../helper}

\title{Introdução à Computação \\ merge sort}
\author{Alexandre Rademaker}
\date{}

\begin{document}

\begin{frame}
 \maketitle
\end{frame}


\begin{frame}[fragile,allowframebreaks]{Ordenação (sort)}

  Dado uma lista de valores, possívelmente desordenada, queremos
  escrever um algorítmo que ordena esta lista, segundo algum critério.

  Por exemplo, para a lista
  \[
    [ 6 \; 3 \; 8 \; 5 \; 2 \; 7 \; 4 \; 1 ]
  \]
  queremos produzir
  \[
    [ 1 \; 2 \; 3 \; 4 \; 5 \; 6 \; 7 \; 8 ]
  \]

  \framebreak

  Vamos considerar dois algorítmos para esta tarefa:

  \begin{itemize}
  \item selection sort
  \item bubble sort
  \item \emph{merge sort}
  \end{itemize}

\end{frame}


\begin{frame}[fragile,allowframebreaks]{Merge sort}

  O \emph{merge sort} é um algorítmo recursivo para ordenação. 

  Composto de duas etapas: \emph{divisão} e \emph{merge}. Na divisão,
  o array inicial é dividido em duas metades e o algorítmo se chama
  para ordenar as metades.

  No merge, dois arrays ordenados são linearmente juntados em um único
  array também ordenado.

  \framebreak

  \begin{lstlisting}[style=code]
    function merge_sort(a[1...n])
      if n > 1:
       x = merge_sort(a[1...n/2])
       y = merge_sort(a[n/2+1...n])
       return merge(x,y)
      else:
       return a

    function merge(x[1...k],y[1...l])
     if k == 0: return y[1...l]
     if l == 0: return x[1...k]
     if x[1] < y[1]:
      return x[1] . merge(x[2...k,y[1...l])
     else:
      return y[1] . merge(x[1...k,y[2...l])
  \end{lstlisting}

  \framebreak

  \begin{center}
    \includegraphics[width=.7\textwidth]{sort-merge.png}
  \end{center}

  \framebreak

  Juntar duas listas ordenadas pode ser feito em tempo $O(n)$ onde $n$
  é o tamanho da lista resultante.

  A quantidade de passos é definida de forma recursiva, $n$ elementos
  geram duas ordenações de $n/2$ e mais um custo linear para o
  \emph{merge}.
  
  \[
    T(n) = 2T(\frac{n}{2}) + O(n)
  \]

  Complexidade é $O(n \log n)$. Por que? Pelo
  \href{https://en.wikipedia.org/wiki/Master_theorem_(analysis_of_algorithms)}{teorema
      mestre}.
      
\end{frame}


\begin{frame}[fragile,allowframebreaks]{A função \texttt{merge}}

  \begin{lstlisting}[style=code]
    2 4 5 7 | 0 1 3 6

    2 4 5 7 | 0 1 3 6
    ^         ^

    2 4 5 7 |   1 3 6
    ^           ^
    0

    2 4 5 7 |     3 6
    ^             ^
    0 1

    4 5 7 |     3 6
    ^           ^
    0 1 2
  \end{lstlisting}

  \framebreak

  \begin{lstlisting}[style=normalc,caption=src/merge.c]
    void merge (int ns[], int l, int m, int u) {
      int t[u - l + 1];
      for (int k = 0, a = l, b = m + 1;
           k < (u - l + 1);
           k++) {
        t[k] =         a > m ? ns[b++] // empty left
             :         b > u ? ns[a++] // empty right
             : ns[a] < ns[b] ? ns[a++] // take left 
             : ns[b++];                // take right
      }
      
      for (int i = l; i <= u; i++)
       ns[i] = t[i - l];
    }
  \end{lstlisting}

\end{frame}


\begin{frame}[fragile,allowframebreaks]{A função \texttt{merge\_sort}}

  \begin{lstlisting}[style=code]
    5 2 7 4 | 1 6 3 0

    5 2 | 7 4

    5 | 2

    (podemos juntar)

    2 5 
  \end{lstlisting}

  Recursivamente divide o array passado em duas metades.

  \framebreak

  \begin{lstlisting}[style=normalc,caption=src/merge.c]
    void merge_sort (int a[], int i, int j) {
      if ((j - i) == 0)
        return;
      
      int m = (i + j) / 2;
      merge_sort(a, i, m);
      merge_sort(a, m + 1, j);  
      merge(a, i, m, j);
    }
  \end{lstlisting}

  \framebreak

  O que fizemos diferente do pseudo-código?

  Ao invés de `retornar' arrays, alteramos os arrays passados.

  O \emph{merge} poderia ser feito `in place', isto é, sem uso de
  memória extra? Se array de entrada tivesse $10^{200}$ posições?

\end{frame}


\begin{frame}[fragile,allowframebreaks]{Conclusão}

  Existem outros algorítmos de ordenação.

  Um site interessante de algorítmos e implementações em várias
  linguagens de programação é o \url{http://rosettacode.org}.

  Podemos
  \href{https://www.cs.usfca.edu/~galles/visualization/ComparisonSort.html}{visualizar}
  e comparar o desempenho dos algorítmos.

  Linguages de programação variam em diversos aspectos, como o foco em
  abstrações de alto nível versus controle preciso de recursos.

  \framebreak

  \begin{lstlisting}[style=code,language=Haskell,basicstyle=\ttfamily\tiny]
  merge :: Ord a => [a] -> [a] -> [a]
  merge [] ys = ys
  merge xs [] = xs
  merge xs@(x:xt) ys@(y:yt)
   | x <= y    = x : merge xt ys
   | otherwise = y : merge xs yt

  split :: [a] -> ([a], [a])
  split []  = ([],[])      
  split [x] = ([x],[])
  split (x:y:zs) = (x:xs,y:ys)
   where (xs,ys) = split zs

  -- >>> split [1,2,3]
  -- ([1,3],[2])

  mergeSort :: Ord a => [a] -> [a]
  mergeSort []  = []
  mergeSort [x] = [x]
  mergeSort xs  = merge (mergeSort as) (mergeSort bs)
   where (as,bs) = split xs

  -- >>> mergeSort [1,2,4,3]
  -- [1,2,3,4]
  \end{lstlisting}

  \framebreak

  Programadores \href{http://haskell.org}{Haskell} adoram abstração (e
  tipos, funções, variáveis imutáveis, pattern matching \ldots).
   
  Mas o código Haskell faz exatamente a mesma coisa que o código C?

  Programadores C adoram controle preciso de recursos!

  \framebreak

  \begin{quote}
    ``The real problem is that programmers have spent far too much
    time worrying about efficiency in the wrong places and at the
    wrong times; premature optimization is the root of all evil (or
    at least most of it) in programming.''
  \end{quote}
  \flushright{\emph{The Art of Computer Programming}, Donald Knuth}

\end{frame}


\end{document}

%%% Local Variables:
%%% mode: latex
%%% TeX-master: t
%%% End:
